% ---- pg-079 （书页 67） ----

\mvsub{矩阵}

\mvbold{定义 2A.13} \quad 一个 $m \times k$ 矩阵是 $m$ 行 $k$ 列个元素组成的矩阵阵列,用黑体大写字母表示,如 $\bm{A},\bm{R},\bm{\Sigma}$ 等.

下面给出几个矩阵实例.

\[
\bm{A} = \begin{bmatrix} -7 & 2 \\ 0 & 1 \\ 3 & 4 \end{bmatrix}, \quad
\bm{B} = \begin{bmatrix} x & 3 & 0 \\ 4 & -2 & 1/x \end{bmatrix}, \quad
\bm{I} = \begin{bmatrix} 1 & 0 & 0 \\ 0 & 1 & 0 \\ 0 & 0 & 1 \end{bmatrix},
\]

\[
\bm{\Sigma} = \begin{bmatrix} 1 & 0.7 & -0.3 \\ 0.7 & 2 & 1 \\ -0.3 & 1 & 8 \end{bmatrix}, \quad
\bm{E} = [\bm{e}_1]
\]

在本书中,矩阵元素是实数或取值为实数的函数.

\mvbold{定义 2A.14} \quad 一个 $m \times k$ 矩阵的维数(缩写为 dim)是有序对 $(m,k)$;$m$ 是行维,$k$ 是列维.通常一个矩阵的维数记在代表矩阵的字母下方的括号中.如 $m \times k$ 矩阵 $\bm{A}$ 用 $\underset{(m \times k)}{\bm{A}}$ 表示.在上述的例子中,矩阵 $\bm{\Sigma}$ 的维数是 $3 \times 3$,可以表示成 $\underset{(3 \times 3)}{\bm{\Sigma}}$.

一个任意常数组成的 $m \times k$ 矩阵 $\bm{A}$,可写做

\[
\underset{(m \times k)}{\bm{A}} = \begin{bmatrix}
a_{11} & a_{12} & \cdots & a_{1k} \\
a_{21} & a_{22} & \cdots & a_{2k} \\
\vdots & \vdots & \ddots & \vdots \\
a_{m1} & a_{m2} & \cdots & a_{mk}
\end{bmatrix}
\]

或更简洁地记做 $\underset{(m \times k)}{\bm{A}} = (a_{ij})$,其中符号 $i$ 代表行,符号 $j$ 代表列.

一个 $m \times 1$ 矩阵叫做列向量,$1 \times k$ 矩阵叫做行向量.由于矩阵可以看做由向量排列组成,因此可以自然地定义数乘和两个同维数的矩阵的加法.

\mvbold{定义 2A.15} \quad 如果 $a_{ij} = b_{ij} ,i=1,2,\cdots ,m,j=1,2,\cdots ,k$,则两个矩阵 $\underset{(m \times k)}{\bm{A}} = (a_{ij})$ 和 $\underset{(m \times k)}{\bm{B}} = (b_{ij})$ 相等,记做 $\bm{A} = \bm{B}$.即如果:

(a) 它们的维数相同,

(b) 各对应元素相同,

则称两矩阵相等.

\mvbold{定义 2A.16} \quad （矩阵加法) \quad 设矩阵 $\bm{A}$ 和 $\bm{B}$ 是两个 $m \times k$ 维的矩阵,其任意元素分别为 $a_{ij}$ 和 $b_{ij} ,i=1,2,\cdots ,m,j=1,2,\cdots ,k$,矩阵 $\bm{A}$ 与 $\bm{B}$ 的和是一个 $m \times k$ 矩阵 $\bm{C}$,记做 $\bm{C} = \bm{A} + \bm{B}$,$\bm{C}$ 的任意元素满足

\[
c_{ij} = a_{ij} + b_{ij} \qquad i=1,2,\cdots ,m,j=1,2,\cdots ,k
\]

注意矩阵加法仅对同维数的矩阵定义.

例如,

\[
\begin{bmatrix} 3 & 2 & 3 \\ 4 & 1 & 1 \end{bmatrix} + \begin{bmatrix} 3 & 6 & 7 \\ 2 & -1 & 0 \end{bmatrix} = \begin{bmatrix} 6 & 8 & 10 \\ 6 & 0 & 1 \end{bmatrix}
\]

\[
\underset{}{\bm{A}} \qquad + \qquad \underset{}{\bm{B}} \qquad = \qquad \underset{}{\bm{C}}
\]

\mvbold{定义 2A.17} \quad （标量乘法) \quad 设 $c$ 为任一标量,$\underset{(m \times k)}{\bm{A}} = (a_{ij})$,则 $c\bm{A} = \bm{A}c = \underset{(m \times k)}{\bm{B}} = (b_{ij})$,其中 $b_{ij} = c a_{ij} = a_{ij} c ,i=1,2,\cdots ,m,j=1,2,\cdots ,k$.
